\section{Computational design}\label{sec:methods}

\subsection{Physical instances and controlled variations}
The base instance contains three vessels per formation at arc-length offsets
$0,2,4$, a known initial geometry, constant operating speed, and bounded
turn controls. The motion polyline has six substeps per epoch. The per-epoch
interaction in \eqref{eq:payoff} is evaluated at endpoints, with no damage
updates or stochastic firing realization. The opponent speed is fixed at
six model distance units per epoch; the focal speed is multiplied by the
declared ratio. Both range ratios equal one. Speed is an imposed operating
condition, not a nested set of selectable speed capabilities.

The directional kernel is the frozen CA-labeled fit. Its target-speed
input remains fixed: a supplied value of six is clipped by the existing
wrapper to five. Accordingly, a motion-speed comparison changes motion
but does not represent a simultaneous recalibration of firing response.
The inherited sector-boundary completion snaps angular differences within
$10^{-8}$ degrees to the specified boundaries. It is applied to the
physical scalar/vector evaluation, not to a projection of the payoff matrix
onto a symmetric matrix.

Grid-3 has normalized controls $\{-1,0,1\}$ multiplied by a maximum turn
of $60$ degrees per epoch. Grid-5 adds $\pm1/2$; Grid-7 uses increments
of $1/3$. Both include Grid-3, but Grid-5 and Grid-7 are not nested in
one another. The horizon remains fixed when calendars are compared.

\begin{table}[t]
\centering\small
\caption{Predetermined computational design. Each physical configuration is evaluated separately against committed and flexible opponents. Shared baseline configurations are counted once in the 45-instance manifest.}\label{tab:design}
\begin{tabular}{p{.19\linewidth}p{.69\linewidth}}
\toprule
Stage & Configuration \\
\midrule
Development & Three geometries; focal speed ratios $0.85,1,1.15$; distance16; $T=6$, Grid-3 \\
Frozen test & Three geometries; paired speed/distance conditions $(.80,13.5),(.95,15),(1.05,17),(1.20,18.5)$; $T=6$, Grid-3 \\
Horizon & Three geometries; speed ratio1, distance16; $T=4,5,6,7$, Grid-3 \\
Action grid & Same fixed anchors; $T=4$, Grid-5 and Grid-7 \\
Ablations & Same anchors at $T=4$, Grid-3; one change at a time: angular averaging, terminal payoff, or rigid line-ahead formation \\
\bottomrule
\end{tabular}
\end{table}
The frozen tests provide 12 physical instances and 24 opponent conditions,
with all 32 calendars at $T=6$. Their speed/distance pairing is a computational
test set, not a factorial experiment identifying the independent effect of
speed. Discovery, implementation amendments and test results retain separate
provenance. Older failed range and mobility interaction hypotheses are not
reclassified as new evidence for the budget problem.

\subsection{Controls, ablations and numerical precision}
Every physical configuration has an independently generated player-swapped
payoff matrix. Parallel geometry additionally requires reflection of signed
turns. The designated speed-one development and first frozen-test anchors
in each geometry also have independently solved mirror calendars. Algebraic
role exchange inside a solver does not substitute for these physical checks.

Angular averaging integrates independent uniform shooter and target
orientations. Shooter-sector weights are $1/6,1/3,1/3,1/6$ and the target
bow/stern weight is $1/3$. With identical kernels this isotropic net
interaction cancels pairwise; it is an analytical control, not an empirical
discovery. The terminal ablation uses $T L(h_T)$ to retain the integrated
objective's time scaling. The rigid line-ahead ablation rotates the whole
formation with its leader instead of following the executed track.

The target numerical gap is $10^{-6}\max(1,\|A\|_\infty)$; response
padding and propagated subgame error are recorded separately. Test cases
include independent micro-game pure-policy enumeration, arbitrary calendars,
the legacy F/C endpoints, full-response reconstruction, information-boundary
checks and non-submodular examples. Deterministic numerical intervals are
not sampling confidence intervals. No population significance test is
attached to a fixed grid of game instances.

\subsection{Comparators and resource accounting}
The comparators are complete calendar enumeration, basic and response-bound
branch-and-bound, exact-increment greedy, and uniform, front-loaded and
back-loaded calendars. The uniform calendar is
$\{\lfloor jT/(K+1)\rfloor:j=1,\ldots,K\}$.
All comparators hold physical input, budget and precision fixed. We report
absolute and interval-valued regret, retained adaptation, full-space response
gaps, oracle calls, runtime and peak memory, including non-improvements.

Two timing measurements are kept distinct. Online trace replay charges
each requested calendar its individually recorded solver cost and executes
response bounds; it measures computational accounting, not end-to-end
wall-clock speedup. Cold runs recompute the prescribed development/test
anchors with an empty method cache and a shared physical payoff input.
Payoff construction is listed separately. Calendar results are reused across
budgets only within the same method.

The experiment uses one heavy local process on a 16GB machine, with an 8GB
process-tree memory limit. Ordinary solver tasks are limited to 20minutes;
bounded continuations use a cumulative six-hour instance budget. Resource
limits, numerical stalls and unresolved values remain in the result record.
No incomplete task is reported as exact convergence.
